\documentclass[12pt]{article}
\usepackage[utf8]{inputenc}
\usepackage[T1]{fontenc}
\usepackage{hyperref}
\title{Yautepec Zapotec Language Resources for K-12 Curriculum Development}
\author{Compiled by Claude Sonnet 5.5 (claude-sonnet-5-5) for Andr\'as Kornai}
\date{}
\begin{document}
\maketitle

\section{Introduction}
Yautepec Zapotec (ISO 639-3 zpb; Glottocode yaut1235) is an endangered Zapotecan language of Oaxaca, Mexico. It is also called Zapoteco de San Bartolo Yautepec or Northwestern Yautepec Zapotec. One mirror of the Endangered Languages Project entry gives about 310 speakers and says the language is unwritten. The search found only catalogue entries and no educational materials.

\section{Primary Educational Resources}
None found.

\section{Linguistic and Reference Materials}
\begin{itemize}
\item Glottolog entry yaut1235 (its bibliography is the best place to look for academic work): \url{https://glottolog.org/resource/languoid/id/yaut1235}
\item Endangered Languages Project, Yautepec Zapotec: \url{https://endangeredlanguages.com/lang/3128}
\item Wikitongues language page: \url{https://wikitongues.org/languages/zpb/}
\end{itemize}

\section{Religious and Community Resources}
Joshua Project reports no Bible portions or other resources: \url{https://joshuaproject.net/languages/zpb}

\section{Digital and Audio Resources}
Global Recordings Network lists no recordings in this language: \url{https://www.globalrecordings.net/en/language/zpb}

\section{Curriculum Development Context}
The language is spoken as a first language by adults only, and no teaching materials, writing system or school use were found. Any curriculum would start with documentation and community-led orthography work. The Glottolog bibliography may point to linguistic descriptions.

\end{document}